\documentclass[UTF8,a4paper,12pt]{ctexart}

\usepackage[a4paper,margin=2.2cm]{geometry}
\usepackage{array}
\usepackage{booktabs}
\usepackage{caption}
\usepackage{enumitem}
\usepackage{float}
\usepackage{fontspec}
\usepackage{hyperref}
\usepackage{listings}
\usepackage{longtable}
\usepackage{pgfplots}
\usepackage{siunitx}
\usepackage{tabularx}
\usepackage{threeparttable}
\usepackage{xcolor}

\pgfplotsset{compat=1.18}
\hypersetup{colorlinks=true,linkcolor=blue!60!black,urlcolor=blue!60!black}
\defaultCJKfontfeatures{AutoFakeSlant=.2}
\setmonofont{Menlo}
\setlength{\parindent}{2em}
\setlength{\parskip}{0.35em}

\definecolor{AccentA}{HTML}{0B6E4F}
\definecolor{AccentB}{HTML}{1D4E89}
\definecolor{AccentC}{HTML}{F4A259}
\definecolor{AccentD}{HTML}{C73E1D}
\definecolor{AccentE}{HTML}{7A306C}
\definecolor{CodeBg}{HTML}{F7F7F7}

\lstset{
  basicstyle=\ttfamily\small,
  backgroundcolor=\color{CodeBg},
  frame=single,
  framerule=0pt,
  xleftmargin=1em,
  xrightmargin=1em,
  breaklines=true,
  columns=fullflexible,
  keywordstyle=\color{AccentB}\bfseries,
  commentstyle=\color{AccentA},
  stringstyle=\color{AccentD},
  showstringspaces=false
}

\title{\heiti 中山大学计算机院本科生实验报告}
\author{}
\date{}

\begin{document}

\begin{titlepage}
  \centering
  {\zihao{2}\bfseries 中山大学计算机院本科生实验报告\par}
  \vspace{0.5em}
  {\zihao{4}（2025学年春季学期）\par}
  \vspace{2em}

  \begin{tabularx}{0.9\textwidth}{>{\bfseries}p{4cm}X}
    课程名称： & 并行程序设计与算法 \\
    实验题目： & Lab0 环境设置与串行矩阵乘法 \\
    批改人： & \\
    专业（方向）： & \\
    学号： & 23336128 \\
    姓名： & 梁力航 \\
    Email： & \\
    完成日期： & 2026年4月1日 \\
  \end{tabularx}

  \vfill
  \begin{minipage}{0.88\textwidth}
    \small
    本报告参考课程实验模板组织内容，围绕串行矩阵乘法的多版本实现、性能测试、结果可视化与优化分析展开。
  \end{minipage}
\end{titlepage}

\section{实验目的}

本实验的目标是实现一个可复现实验流程的串行矩阵乘法程序，并通过多版本对照分析不同实现策略对性能的影响。具体目标如下：

\begin{enumerate}[leftmargin=2.2em]
  \item 使用 \texttt{Python} 与 \texttt{C++} 分别实现基础串行矩阵乘法；
  \item 比较循环顺序、编译优化和循环展开对缓存局部性及运行时间的影响；
  \item 用统一随机种子、统一校验值和统一统计口径保证结果可比性；
  \item 以表格和图表展示实验结果，并分析各版本的收益边界；
  \item 补充 \texttt{Intel MKL} 版本在 Docker 异构环境中的运行路径，并说明该数据仅用于功能验证与接口完备性说明。
\end{enumerate}

\section{实验过程和核心代码}

\subsection{实验平台与实现结构}

本次实验在 \texttt{macOS arm64} 平台上完成，编译器为 \texttt{Apple clang 17.0.0}，解释器为 \texttt{Python 3.12.4}。为保证不同版本具备可比性，程序统一采用以下约束：

\begin{itemize}[leftmargin=2.2em]
  \item 统一输入协议：\texttt{m n k [seed]}；
  \item 统一矩阵生成方式：基于相同随机种子生成双精度矩阵；
  \item 统一结果校验：输出矩阵校验和 \texttt{checksum}；
  \item 统一性能指标：运行时间、相对加速比、绝对加速比与 GFLOPS。
\end{itemize}

本实验共实现 6 个版本：

\begin{enumerate}[leftmargin=2.2em]
  \item \textbf{Python}：纯 Python 三重循环，作为绝对基线；
  \item \textbf{C++基础版}：最直接的 \texttt{i-j-p} 三重循环；
  \item \textbf{调整循环顺序}：采用 \texttt{i-p-j} 顺序，提升缓存友好性；
  \item \textbf{编译优化}：在版本 3 基础上启用 \texttt{-O3}；
  \item \textbf{循环展开}：在版本 4 基础上进行 4 路展开；
  \item \textbf{Intel MKL}：使用 \texttt{cblas\_dgemm}，通过 Docker 的 \texttt{linux/amd64} oneMKL 环境获得实测结果。
\end{enumerate}

\subsection{核心代码片段}

\subsubsection{基础版本：\texttt{i-j-p} 三重循环}

\begin{lstlisting}[language=C++]
for (int i = 0; i < m; ++i) {
    for (int j = 0; j < k; ++j) {
        double sum = 0.0;
        for (int p = 0; p < n; ++p) {
            sum += a[i * n + p] * b[p * k + j];
        }
        c[i * k + j] = sum;
    }
}
\end{lstlisting}

该版本逻辑最直观，但访问矩阵 \texttt{B} 时存在明显的跨行跳跃，缓存局部性较差。

\subsubsection{循环重排：\texttt{i-p-j}}

\begin{lstlisting}[language=C++]
for (int i = 0; i < m; ++i) {
    for (int p = 0; p < n; ++p) {
        const double a_ip = a[i * n + p];
        for (int j = 0; j < k; ++j) {
            c[i * k + j] += a_ip * b[p * k + j];
        }
    }
}
\end{lstlisting}

经过循环重排后，\texttt{B} 和 \texttt{C} 的访问模式更接近顺序扫描，因此缓存命中率显著提高。

\subsubsection{循环展开}

\begin{lstlisting}[language=C++]
for (int j = 0; j + 3 < k; j += 4) {
    c[c_row + j]     += a_ip * b[b_row + j];
    c[c_row + j + 1] += a_ip * b[b_row + j + 1];
    c[c_row + j + 2] += a_ip * b[b_row + j + 2];
    c[c_row + j + 3] += a_ip * b[b_row + j + 3];
}
\end{lstlisting}

循环展开可以减少分支判断与循环控制开销，但在现代编译器已做较强自动优化时，额外收益通常不会像 \texttt{-O3} 那样显著。

\subsection{正确性验证方法}

本实验采用固定随机种子 \texttt{20250401} 生成输入矩阵。不同版本在同一输入上的输出校验值应保持一致。以主实验规模 \texttt{512} $\times$ \texttt{512} $\times$ \texttt{512} 为例，宿主机原生版本 \texttt{v1-v5} 的校验值统一为：

\begin{center}
\fbox{\parbox{0.82\textwidth}{\centering
\texttt{checksum = -6740.486084331}
}}
\end{center}

这说明宿主机原生的版本 1 至版本 5 在数值上保持一致，性能对比具有有效性。版本 6 在 Docker 的 \texttt{linux/amd64} 异构环境下得到 \texttt{checksum = -6740.486084332}，与原生结果在当前输出精度和比较容差下保持一致，可用于确认实现正确，但不参与严格公平的原生性能比较。

\section{实验结果}

\subsection{主实验结果汇总}

主实验采用矩阵规模 \texttt{512} $\times$ \texttt{512} $\times$ \texttt{512}，每个版本重复执行 3 次，取平均时间。由于当前设备峰值浮点性能未进行严格标定，表中的“峰值性能百分比”统一记为 \texttt{N/A}，避免伪精确。

\begin{table}[H]
  \centering
  \caption{主实验结果汇总（512\,$\times$\,512\,$\times$\,512）}
  \begin{threeparttable}
    \small
    \renewcommand{\arraystretch}{1.2}
    \begin{tabularx}{\textwidth}{
      >{\centering\arraybackslash}m{0.8cm}
      >{\raggedright\arraybackslash}X
      >{\centering\arraybackslash}m{1.55cm}
      >{\centering\arraybackslash}m{1.65cm}
      >{\centering\arraybackslash}m{1.65cm}
      >{\centering\arraybackslash}m{1.4cm}
      >{\centering\arraybackslash}m{1.1cm}
    }
      \toprule
      版本 & 实现描述 & \shortstack{运行时间\\(s)} & \shortstack{相对\\加速比} & \shortstack{绝对\\加速比} & \shortstack{GF\\LOPS} & \shortstack{峰值\\\%} \\
      \midrule
      1 & Python & 4.7499 & 1.0000 & 1.0000 & 0.0565 & N/A \\
      2 & C++基础版 & 0.4252 & 11.1723 & 11.1723 & 0.6314 & N/A \\
      3 & 调整循环顺序 & 0.1533 & 2.7738 & 30.9894 & 1.7513 & N/A \\
      4 & 编译优化（-O3） & 0.0178 & 8.6290 & 267.4082 & 15.1123 & N/A \\
      5 & 循环展开 & 0.0177 & 1.0055 & 268.8763 & 15.1952 & N/A \\
      6 & Intel MKL（Docker/x86\_64）\tnote{a} & 0.2656 & 0.0665 & 17.8829 & 1.0106 & N/A \\
      \bottomrule
    \end{tabularx}
    \begin{tablenotes}[flushleft]
      \footnotesize
      \item[a] 当前机器为 Apple Silicon（Apple M4 / macOS arm64）。Intel 官方 macOS oneMKL 路线面向 Intel 64，且 oneAPI 2024.0 起不再支持 macOS。因此，版本 6 通过 Docker 的 \texttt{linux/amd64} 仿真环境运行 x86\_64 Linux 版 oneMKL 得到。该次实测在 \texttt{512}$\times$\texttt{512}$\times$\texttt{512} 下耗时 \num{0.2656} 秒，校验值为 \texttt{-6740.486084332}，仅用于功能完成与接口验证，不与本机原生版本 1--5 做严格公平的性能比较。
    \end{tablenotes}
  \end{threeparttable}
\end{table}

\subsection{可视化结果}

由于 \texttt{v6\_mkl} 来自 Docker 的 \texttt{linux/amd64} 异构环境，为避免混入口径不一致的数据，以下图表仅展示宿主机原生版本的可比结果。

\begin{figure}[H]
  \centering
  \begin{tikzpicture}
    \begin{axis}[
      width=0.94\textwidth,
      height=7.2cm,
      ybar,
      bar width=14pt,
      ylabel={运行时间（秒，对数坐标）},
      xlabel={版本},
      ymode=log,
      log basis y=10,
      symbolic x coords={Python,C++基础版,调整循环顺序,编译优化,循环展开},
      xtick=data,
      x tick label style={rotate=15,anchor=east,font=\small},
      ymin=0.01,
      ymax=10,
      nodes near coords,
      nodes near coords align={vertical},
      every node near coord/.append style={font=\scriptsize, rotate=90, anchor=west},
      enlarge x limits=0.08,
      grid=both,
      minor grid style={draw=gray!15},
      major grid style={draw=gray!30}
    ]
      \addplot[fill=AccentB!75, draw=AccentB!90] coordinates {
        (Python,4.7499)
        (C++基础版,0.4252)
        (调整循环顺序,0.1533)
        (编译优化,0.0178)
        (循环展开,0.0177)
      };
    \end{axis}
  \end{tikzpicture}
  \caption{各版本运行时间对比（对数坐标）}
\end{figure}

\begin{figure}[H]
  \centering
  \begin{tikzpicture}
    \begin{axis}[
      width=0.94\textwidth,
      height=7.2cm,
      ybar,
      bar width=14pt,
      ylabel={相对 Python 的绝对加速比（对数坐标）},
      xlabel={版本},
      ymode=log,
      log basis y=10,
      symbolic x coords={Python,C++基础版,调整循环顺序,编译优化,循环展开},
      xtick=data,
      x tick label style={rotate=15,anchor=east,font=\small},
      ymin=1,
      ymax=400,
      nodes near coords,
      nodes near coords align={vertical},
      every node near coord/.append style={font=\scriptsize, rotate=90, anchor=west},
      enlarge x limits=0.08,
      grid=both,
      minor grid style={draw=gray!15},
      major grid style={draw=gray!30}
    ]
      \addplot[fill=AccentC!80, draw=AccentD!90] coordinates {
        (Python,1.0000)
        (C++基础版,11.1723)
        (调整循环顺序,30.9894)
        (编译优化,267.4082)
        (循环展开,268.8763)
      };
    \end{axis}
  \end{tikzpicture}
  \caption{各版本相对 Python 的绝对加速比}
\end{figure}

\begin{figure}[H]
  \centering
  \begin{tikzpicture}
    \begin{axis}[
      width=0.94\textwidth,
      height=7.6cm,
      xlabel={矩阵规模（方阵边长）},
      ylabel={GFLOPS},
      xmin=480, xmax=1056,
      ymin=0, ymax=17,
      xtick={512,768,1024},
      legend style={at={(0.5,1.02)},anchor=south,legend columns=4,draw=none,font=\small},
      grid=major,
      line width=1pt,
      mark size=2.8pt
    ]
      \addplot[color=AccentB,mark=*] coordinates {(512,0.6256) (768,0.9983) (1024,0.5782)};
      \addlegendentry{C++基础版}
      \addplot[color=AccentA,mark=square*] coordinates {(512,1.7632) (768,1.7645) (1024,1.7703)};
      \addlegendentry{调整循环顺序}
      \addplot[color=AccentD,mark=triangle*] coordinates {(512,15.0539) (768,14.9698) (1024,15.0365)};
      \addlegendentry{编译优化}
      \addplot[color=AccentE,mark=diamond*] coordinates {(512,14.9927) (768,14.9681) (1024,14.9858)};
      \addlegendentry{循环展开}
    \end{axis}
  \end{tikzpicture}
  \caption{C++ 各优化版本在不同规模下的 GFLOPS 变化}
\end{figure}

\subsection{结果分析}

\begin{enumerate}[leftmargin=2.2em]
  \item \textbf{Python 与 C++ 的差距极大。} 在 $512^3$ 规模下，Python 版本平均耗时约 \SI{4.75}{\second}，而最直接的 C++ 实现仅需 \SI{0.43}{\second}，绝对加速比达到约 \num{11.17}。这说明解释器开销对数值密集型三重循环影响非常明显。

  \item \textbf{循环顺序优化带来了稳定而显著的收益。} 从版本 2 到版本 3，运行时间由 \SI{0.4252}{\second} 降至 \SI{0.1533}{\second}，相对加速比约为 \num{2.77}。其根本原因是 \texttt{i-p-j} 顺序使矩阵 \texttt{B} 和 \texttt{C} 的访问更连续，显著改善缓存局部性。

  \item \textbf{编译优化是收益最大的单次改动。} 从版本 3 到版本 4，时间进一步降至 \SI{0.0178}{\second}，相对加速比达到 \num{8.63}，GFLOPS 从 \num{1.75} 提升到 \num{15.11}。这表明编译器在指令级优化、寄存器分配和循环层面优化方面的能力非常强。

  \item \textbf{手工循环展开的边际收益有限。} 版本 5 相对版本 4 的加速比仅为 \num{1.0055}，说明在当前编译器和硬件环境下，\texttt{-O3} 已经完成了大部分关键优化，继续手工展开只能带来非常有限的收益。

  \item \textbf{优化版本的 GFLOPS 随规模变化较稳定。} 从图 3 可见，版本 4 和版本 5 在 \texttt{512}、\texttt{768}、\texttt{1024} 三组规模下均维持在约 \num{15} GFLOPS，说明该实现已较充分利用当前平台的串行浮点能力；而未优化版本明显更受访存模式和指令效率制约。

  \item \textbf{MKL 版本已在 Docker 异构环境中跑通，但不能按原生口径比较。} 当前机器为 \texttt{Apple M4 / macOS arm64}。Intel 官方 macOS oneMKL 路线面向 \texttt{Intel 64}，且从 \texttt{oneAPI 2024.0} 起已停止 macOS 支持。因此，版本 6 只能通过 Docker 的 \texttt{linux/amd64} 仿真容器运行 Linux 版 oneMKL。本实验已在该环境下得到正确结果，\texttt{checksum} 在当前输出精度和比较容差下与原生版本保持一致，但由于运行环境异构，表中的时间、GFLOPS 和加速比只表示“在仿真环境中能完成计算”，不应与宿主机原生版本 1--5 做严格公平的性能比较。
\end{enumerate}

\section{实验感想}

本次实验最直接的体会是：矩阵乘法这类数值计算问题看似“算法固定”，但工程实现细节会带来数量级上的性能差异。仅仅把 Python 切换到 C++，性能就获得了十倍以上提升；而在 C++ 中进一步调整循环顺序，又能仅凭缓存友好性获得接近三倍收益。真正将性能推到一个更高水平的，是编译器优化与底层代码生成能力。

另一个值得注意的结论是，手工优化并不总是越多越好。在本实验环境中，循环展开只比 \texttt{-O3} 略快，说明现代编译器已经能够自动完成大量低层优化。相比盲目堆叠技巧，更重要的是先定位瓶颈，再做有针对性的优化。

最后，性能实验必须建立在正确性一致的基础上。本次实验通过统一随机种子与校验和比较，确保宿主机原生版本在数值上保持一致，并确认 \texttt{v6\_mkl} 在 Docker 异构环境下也通过了 \texttt{checksum} 校验。只有在正确结果的前提下，性能曲线和加速比才有分析价值。

\end{document}
